% Part: applied-modal-logics
% Chapter: temporal-logic
% Section: properties-accessibility

\documentclass[../../../include/open-logic-section]{subfiles}

\begin{document}

\olfileid[id]{aml}{tl}{acc}

\olsection{Sifat Kerangka Temporal}

Karena model temporal kita tidak mengenakan syarat apa pun pada
relasi~$\prec$, satu-satunya aksioma yang telah kita kenal dan berlaku dalam
semua model ialah~$K$, atau analognya $K_{\Gtemp}$ dan~$K_{\Htemp}$:
\begin{align*}
\tag{$K_{\Gtemp}$}  \Gtemp (p \lif q) & \lif (\Gtemp p \lif \Gtemp q)\\
\tag{$K_{\Htemp}$}  \Htemp (p \lif q) & \lif (\Htemp p \lif \Htemp q)
\end{align*}

Namun, jika kita ingin model kita mengenakan syarat yang lebih ketat pada
cara waktu direpresentasikan, misalnya dengan memastikan bahwa $\prec$
merupakan urutan linear, model kita akan mempunyai validitas lain.

\begin{table}[t]
  \begin{tabular}{| p{.48\textwidth} || p{.48\textwidth} |}
    \hline
    {\emph{Jika $\prec$ bersifat \dots}} & {\emph{maka \dots benar
    dalam~$\mModel{M}$:}} \\
    \hline \hline
    \emph{transitif}: \newline
    $\forall u \forall v \forall w ((u \prec v \land v \prec w) \lif u \prec w)$ &
    $\Ftemp \Ftemp p \lif \Ftemp p$  \\
    \hline
    \emph{linear}: \newline
    $\forall w \forall v (w \prec v \lor w = v \lor v \prec w)$ &
    $(\Ftemp \Ptemp p \lor \Ptemp \Ftemp p) \lif
    (\Ptemp p \lor p \lor \Ftemp p)$\\
    \hline
    \emph{rapat}: \newline
    $\forall w \forall v (w \prec v \lif \exists u(w \prec u \land u \prec v))$ &
    $\Ftemp p \lif \Ftemp \Ftemp p$ \\
    \hline
    \emph{tak berbatas (masa lampau)}: \newline
    $\forall w \exists v( v \prec w)$ &
    $\Htemp p \lif \Ptemp p$ \\
    \hline
    \emph{tak berbatas (masa depan)}: \newline
    $\forall w \exists v( w \prec v)$ &
    $\Gtemp p \lif \Ftemp p$ \\
    \hline
  \end{tabular}
  \caption{Beberapa sifat korespondensi kerangka temporal.}
  \ollabel{tab:correspondence}
\end{table}

Beberapa sifat dalam \olref{tab:correspondence} mungkin tampak sebagai ciri
yang diinginkan dari model yang dimaksudkan untuk merepresentasikan waktu.
Namun, perlu diperhatikan bahwa, walaupun kita dapat mengenakan syarat apa pun
yang kita inginkan pada relasi~$\prec$, tidak semua syarat berkorespondensi
dengan !!{formula}s yang dapat dinyatakan dalam bahasa logika temporal.
Sebagai contoh, irefleksivitas, yakni gagasan bahwa
$\forall w\lnot(w\prec w)$, tidak mempunyai formula yang berkorespondensi
dengannya dalam logika temporal.

\end{document}
